\documentclass[12pt, twoside, italian]{article}
\input{h_preamble/new_preamble}
\input{h_preamble/letterfonts}
\input{h_preamble/macros}
\usepackage{fancyhdr}
\usepackage{titlesec}  % Opzionale, per controllo avanzato sui titoli
\setlength{\parindent}{0pt}
\setlength{\parskip}{0pt}
\pagestyle{fancy}

% Definizione dell'intestazione per pagine dispari (Odd)
\fancyhead[LO]{\textbf{\thepage}}  % Numero di pagina in grassetto a sinistra
\fancyhead[RO]{\nouppercase\leftmark}         % Capitolo e sezione a destra

% Definizione dell'intestazione per pagine pari (Even)
\fancyhead[LE]{\nouppercase\rightmark}          % Capitolo e sezione a sinistra
\fancyhead[RE]{\textbf{\thepage}}  % Numero di pagina in grassetto a destra

% Cancella il footer
\fancyfoot{}
\setlength{\headheight}{14.49998pt}
\addtolength{\topmargin}{-2.49998pt}
% Personalizzazione della linea orizzontale
\renewcommand{\headrulewidth}{0.4pt}
\renewcommand{\footrulewidth}{0pt}

\pgfplotsset{compat=newest}
\title{Simulare buchi neri}
\author{Francesco Sermi}
\date{\today}
\pgfplotsset{
	compat = newest,
	colormap/violet % se si vuole cambiare il colormap utilizzato per generare i grafici, tuttavia la copertina ha un colormap a
					% parte che va modificato qua sotto. Se si stampa non a colori consiglio blackwhite come colormap
	}
\renewcommand{\appendixpage}{}
\begin{document}
	%\begin{titlepage}
	%\centering
	%\vspace*{3cm}
	%{\huge\bfseries Appunti di Geometria e Algebra Lineare \par}
	%\vspace{2cm}
	%{\Large\itshape Francesco Sermi\par}
	%\vfill
	%\begin{figure}[H]
	%	\centering
	%	\includegraphics[width=0.65\textwidth]{immagini/dejong.png}
	%\end{figure}
    	%\begin{tikzpicture}
	% Definisco lo stile dei vettori
    	%\tikzset{vector/.style={-stealth, thick, color=black}}

    % Parametri per il punto di contatto
    	%\def\xo{1}
    	%\def\yo{2}
    	%\def\zo{5} % Calcola z0 = f(x0, y0)

    % Derivate parziali della funzione f(x, y) = x^2 + y^2 nel punto (x0, y0)
    	%\def\fx{2*\xo} % = 2
    	%\def\fy{2*\yo} % = 4

    % Disegno la superficie
    	%\begin{axis}[
    	%    width=12cm,
    	%    view={60}{30}, % Angolazione impostata
    	%    axis lines=middle,
    	%    zlabel={$f(x,y)$},
    	%    clip=false,
    	%    domain=-5:5,
    	%    y domain=-5:5,
    	%    samples=30,
    	%    colormap/viridis,
    	%    z buffer=sort
    	%]

    % Grafico della superficie con colore visibile
    	%\addplot3[surf, opacity=0.8, shader=interp]
    	%    {x^2 - y^2}; % Superficie f(x, y) = x^2 + y^2
    	%\end{axis}
	%\end{tikzpicture}
	%\vfill	
	%\end{titlepage}
	%\pagestyle{empty}
	%\vspace*{\fill}
    %Appunti di Geometria e Algebra Lineare  © 2025 by Francesco Sermi is licensed under Creative Commons Attribution-NonCommercial-NoDerivatives 4.0 International. To view a copy of this license, visit https://creativecommons.org/licenses/by-nc-nd/4.0/
	%\cleardoublepage
	%\vspace*{8cm}
	\maketitle
	
	Lo scopo di questo breve (o almeno spero) documento è quello di riportare in forma abbastanza elegante tutti i risultati matematici e spiegazioni necessarie per simulare un buco nero in Python (o in C++ se vedo che non è possibile). 
	
	\section{Requisiti matematicosi}

	Il punto di partenza è sicuramente la metrica di Schwarzschild
	$$
		ds^2 = -(1 - \frac{2GM}{r})dt^2 + (1 - \frac{2GM}{r})^{-1}dr^2 + r^2(d\theta^2 + \sin^2{\theta}d\phi^2),
	$$
	da cui segue che il nostro tensore metrico
	$$
		g_{\mu \nu} = \begin{pmatrix}
			-(1 - \frac{2GM}{r}) & 0 & 0 & 0 \\
			0 & (1-\frac{2GM}{r})^{-1} & 0 & 0 \\
			0 & 0 & r^2 & 0 \\
			0 & 0 & 0 & r^2 \sin^2{\theta}
		\end{pmatrix}
	$$
	Visto che tengo alla mia sanità mentale, vorrei, innanzitutto, introdurre la notazione di Einstein che mi permetterà di snellire particolarmente le formule: se abbiamo per esempio due vettori, $v$ e $w$, le cui componenti verranno indicate con $v^\mu$ e $w^\mu$ rispettivamente, avremo che scrivere
	$$
		v^\mu w_\mu,
	$$
	dove abbiamo un indice ripetuto indica che abbiamo una sommatora non esplicita su tale indice, dunque
	$$
		v^\mu w_\mu = \sum_{\mu=0}^3 v^\mu w_\mu = v^0 w_0 + v^1 w_1 + v^2 w_2 + v^3 w_3.
	$$
	In accordo con la letteratura, diremo che un oggetto con un indice in alto trasforma in maniera \emph{controvariante} altrimenti trasforma in maniera \emph{covariante}:
	\begin{definition}[trasformazioni covarianti e controvarianti]
		Diremo che un oggetto $v$ che da un sistema di riferimento ad un altro (con jacobiano $\xi^\mu_\nu = \frac{\partial (x')^\mu}{\partial x^\nu}$) trasforma in maniera controvariante se
		$$
			(v')^\mu = \xi^\mu_\nu v^\nu,
		$$
		e gli posizioneremo l'indice in alto. Altrimenti, diremo che trasforma in maniera covariante se
		$$
			(v')_\mu = \xi^\nu_\mu v_\nu,
		$$
		e gli posizioneremo l'indice in basso.
	\end{definition}
	Chiaramente esistono degli oggetti con più indici: i tensori, che possono avere $p$ indici in alto e $q$ indici in basso, che sono delle applicazioni che mappano $p$ vettori e $q$ covettori (funzionali) in un numero\footnote{si possono dare anche altre definizioni (su varietà per esempio), ma ci vogliamo tenere sul semplice.},
	$$
		\Lambda^{\alpha_1, \ldots, \alpha_p}_{\beta_1, \ldots, \beta_q},
	$$
	che trasforma come segue
	$$
		(\Lambda')^{\mu_1, \ldots, \mu_p}_{\nu_1, \ldots, \nu_q} = \xi^{\mu_{1}}_{\alpha_1} \xi^{\mu_{2}}_{\alpha_2} \ldots \xi^{\mu_p}_{\alpha_p} \xi^{\beta_1}_{\nu_1} \xi^{\beta_2}_{\nu_2} \ldots \xi^{\nu_q}_{\beta_q} \Lambda^{\alpha_1, \ldots, \alpha_p}_{\beta_1, \ldots, \beta_q}.
	$$
	Abbiamo che i simboli di Christoffel in questa metrica sono dati da
	\begin{align*}
		&\Gamma^{\sigma}_{\mu \nu} = \frac{1}{2}g^{\sigma \rho}(\partial_{\mu} g_{\nu \rho} + \partial_{\nu} g_{\rho \mu} - \partial_\rho g_{\mu \nu}), &  &\partial_\mu = \frac{\partial}{\partial x^\mu},
	\end{align*}
	Facendo il conto\footnote{si potrebbe anche pensare di fare noi stessi il conto, ma voglio alleggerire il carico di lavoro al mio povero PC, quindi si fa analiticamente il conto} (facendolo fare a Wolfram Mathematica) si vede che
	\begin{align*}
		&\Gamma^\phi_{ij} = \begin{pmatrix}
			0 & 0 & 0 & 0 \\
			0 & 0 & 0 & \frac{1}{r} \\
			0 & 0 & 0 & \cot{\theta} \\
			0 & \frac{1}{r} & \cot{\theta} & 0 
		\end{pmatrix}, & &\Gamma^{\theta}_{ij} = \begin{pmatrix}
			0 & 0 & 0 & 0 \\
			0 & 0 & \frac{1}{r} & 0 \\
			0 & \frac{1}{r} & 0 & 0 \\
			0 & 0 & 0 & -\cos{\theta}\sin{\theta}
		\end{pmatrix}, \\
		&\Gamma^{r}_{ij} = \begin{pmatrix}
			0 & \frac{GM}{r^2}(1-\frac{2GM}{r})^{-1} & 0 & 0 \\
			\frac{GM}{r^2}(1-\frac{2GM}{r})^{-1} & 0 & 0 & 0 \\
			0 & 0 & 0 & 0 \\
			0 & 0 & 0 & 0
		\end{pmatrix}
	\end{align*}
	\begin{definition}[metrica singolare]
		Diremo che una metrica è singolare se $\exists x \in M$ con $M$ varietà riemmaniana tale che $g_{\mu \nu}$ non è definito.
	\end{definition}
	Vediamo che la metrica di Schwarzschild non è ben definita se $r = 2GM$ e per $r = 0$: la quantità $r_S = 2GM$ verrà indicata con \emph{raggio di Schwarzschild}.
	Da Geometria differenziale sappiamo che le geodetiche\footnote{presi due punti $p, q$ distinti su una varietà, esiste un'unica curva (che chiamiamo \emph{geodetica}) di lunghezza minima che collega i due punti.} rispettano la seguente equazione:
	\begin{align}
		&\frac{\partial^2 x^\mu}{\partial t^2} + \frac{\partial x^i}{\partial t} \frac{\partial x^j}{\partial t} \Gamma^k_{ij} = 0.
	\end{align}
	Quindi, il nostro scopo, sarà quello di generare un oggetto molto massivo (er buco nero) e andare ad integrare numericamente queste equazioni per ricostruire il percorso della luce: chiarissimo, no? Detto questo, veniamo alla scelta dell'integratore: visto che ancora non so
	niente perché ancora non ho ancora "messo le mani in pasta" (ovvero non so ancora quanto sia brutta effettivamente l'equazione, che dev'essere ancora discretizzata), effettuerò un bellissimo \emph{educated guess} e userò l'integratore Runge-Kutta. Il primo step che dobbiamo impiegare
	è quello di costruire una funzione che ci calcola il tensore metrico per ogni punto considerado
	\begin{lstlisting}[language=Python]
def schwarzschild_metric(position){
	rs = 1
	r = position[1]
	theta = position[2]
	g = np.zeroes(shape=[4, 4])
	g[0][0] = -(1-rs/r)
	g[1][1] = 1/(1-rs/r)
	g[2][2] = r*r
	g[3][3] = r * r * np.sin(theta) * np.cos(theta)
	return g
}
	\end{lstlisting}
	Successivamente bisogna calcolare le tetradi: localmente, lo spazio curvo deve assomigliare ad uno spazio euclideo, quindi bisogna calcolare i vettori di base dello spazio euclideo "curvati" dalla curvatura dello spazio-tempo: nel caso della metrica di Schwarzschild, abbiamo che questi valori sono già noti in letteratura e risultano essere pari a
	\begin{align}
		&e_0 = (\frac{1}{\sqrt{1-r_s/r}}, 0, 0, 0) \\
		&e_1 = (0, \frac{1}{\sqrt{1-r_s/r}}, 0, 0) \\
		&e_2 = (0, 0, \frac{1}{r}, 0) \\
		&e_3 = (0, 0, 0, \frac{1}{r\sin{\theta}}).
	\end{align}
	\begin{lstlisting}[language=Python]
def schwarzschild_tetrad(position):
	rs = 1
	r = position[1]
	theta = position[2]
	e0 = np.array([1/(np.sqrt(1-(rs/r))), 0, 0, 0])
	e1 = np.array([0, 1/np.sqrt(1-rs/r)], 0, 0)
	e2 = np.array([0, 0, 1/r, 0])
	e3 = np.array([0, 0, 0, 1/(r * np.sin(theta))])
	return np.array([e0, e1, e2, e3])
\end{lstlisting}
	Ora, noi vogliamo vedere tutto questo da una bella camera, dunque dobbiamo convertire i nostri dati proiettandoli su uno schermo: per fare questo si fa un pochino di trigonometria (che vi risparmio, perché geometria euclidea mi fa schifo, bensì la fucking geometria non euclidea mi garba un mondo): si converte il FOV in radianti per poi calcolare la distanza massima rispetto al piano di proiezione (la camera) sapendo che
	$$
		d = \frac{l}{2\tan{\frac{\text{fov}}{2}}},
	$$
	per quanto riguarda le posizioni $(x, y)$, abbiamo che se posizioniamo il nostro \emph{Super massive Black Hole} nell'origine e prendiamo il sistema di riferimento da cui prendere le misure lì allora
	\begin{align*}
		&x = s_x - \frac{l_x}{2}, \\
		&y = s_y - \frac{l_y}{2}.
	\end{align*}
	Una volta che ci siamo calcolati la direzione del vettore, andiamo a normalizzarlo.
\end{document}
Andiamo alle geodetiche: per adesso proviamo ad integrarle direttamente
\begin{lstlisting}
def geodetic(position, direction, tetrads):
	g = {
		"position": [],
		"velocity": [0, 0, 0, 0]
	}
	g = tetrads[0] * -1 + tetrads[1] * direction[0] + tetrads[2] * direction[1] + tetrads[3] * direction[2]
	return g
\end{lstlisting}
Infine si scrive una funzione che mi dà i simbolozzi di Christoffel e mi integra
\begin{lstlisting}
	#-------------------------------------------
# Numerical derivative: central difference
#-------------------------------------------
def diff(func, position, direction):
    h = 1e-5
    p_up = position.copy()
    p_lo = position.copy()

    p_up[direction] += h
    p_lo[direction] -= h

    return (func(p_up) - func(p_lo)) / (2 * h)


#--------------------------------------------------
# Compute Christoffel symbols Γ^μ_{αβ}
# get_metric : function(position) → 4×4 metric
#--------------------------------------------------
def calculate_christoff2(position, get_metric):
    g = get_metric(position)                # g_{μν}
    g_inv = np.linalg.inv(g)                # g^{μν}

    metric_diff = np.zeros((4, 4, 4))       # ∂_i g_{jk}

    # Compute derivatives of metric in all 4 directions
    for i in range(4):
        differentiated = diff(get_metric, position, i)   # 4×4 matrix
        for j in range(4):
            for k in range(4):
                metric_diff[i, j, k] = differentiated[j, k]

    # Christoffel symbols Γ^μ_{αβ}
    Gamma = np.zeros((4, 4, 4))

    for mu in range(4):
        for al in range(4):
            for be in range(4):
                s = 0.0
                for sigma in range(4):
                    s += 0.5 * g_inv[mu, sigma] * (
                        metric_diff[be, sigma, al]
                        + metric_diff[al, sigma, be]
                        - metric_diff[sigma, al, be]
                    )
                Gamma[mu, al, be] = s

    return Gamma


#-------------------------------------------------------
# Geodesic equation: a^μ = - Γ^μ_{αβ} v^α v^β
#-------------------------------------------------------
def calculate_acceleration_of(X, v, get_metric):
    Gamma = calculate_christoff2(X, get_metric)
    a = np.zeros(4)

    for mu in range(4):
        s = 0.0
        for al in range(4):
            for be in range(4):
                s += -Gamma[mu, al, be] * v[al] * v[be]
        a[mu] = s

    return a


#-------------------------------------------------------
# Schwarzschild specialization
#-------------------------------------------------------
def calculate_schwarzschild_acceleration(X, v):
    return calculate_acceleration_of(X, v, schwarzschild_metric)


#-------------------------------------------------------
# Integration result enum
#-------------------------------------------------------
class IntegrationResult:
    ESCAPED = 0
    EVENT_HORIZON = 1
    UNFINISHED = 2

    def __init__(self):
        self.type = IntegrationResult.UNFINISHED


#-------------------------------------------------------
# Leapfrog integrate a geodesic
# g.position and g.velocity = numpy arrays (4,)
#-------------------------------------------------------
def integrate(g, debug=False):
    result = IntegrationResult()

    dt = 0.005
    rs = 1.0
    start_time = g.position[0]

    for _ in range(100000):
        acceleration = calculate_schwarzschild_acceleration(g.position, g.velocity)

        # Leapfrog step
        g.velocity = g.velocity + acceleration * dt
        g.position = g.position + g.velocity * dt

        radius = g.position[1]

        # Escape condition
        if radius > 10.0:
            result.type = IntegrationResult.ESCAPED
            return result

        # Event horizon or timeout
        if radius <= rs + 1e-4 or g.position[0] > start_time + 1000:
            result.type = IntegrationResult.EVENT_HORIZON
            return result

    return result
\end{lstlisting}